\section{Method}
\label{sec:method}

\subsection{Register-augmented vision transformer}

Let $x \in \mathbb{R}^{3 \times H \times W}$ be an input image, partitioned by a
$16 \times 16$ patch embedding into $N = 196$ patch tokens of width $d = 192$.
The standard ViT prepends a classification token, giving the sequence
$[\,\texttt{[CLS]}, p_1, \dots, p_N\,]$. We insert $K$ learnable register
tokens $r_1 \dots r_K \in \mathbb{R}^{d}$ between the two groups:

\begin{equation}
z^{(0)} = [\,\texttt{[CLS]},\; r_1, \dots, r_K,\; p_1, \dots, p_N\,] + E_{\mathrm{pos}},
\end{equation}

where positional embeddings $E_{\mathrm{pos}}$ are applied to the
\texttt{[CLS]} and patch positions of the pretrained backbone; registers carry
no positional embedding, as they correspond to no spatial location. The
sequence length becomes $S = 1 + K + N$. The registers participate fully in
every self-attention block and are discarded before the classification head,
which reads the \texttt{[CLS]} token only. The register block is initialised
with a truncated normal, $r_i \sim \mathcal{N}_{\mathrm{trunc}}(0, 0.02^2)$;
zero initialisation is avoided because it makes the register columns
indistinguishable and stalls their gradient.

The parameter cost is $K \cdot d$ weights, which is negligible: $K=8$ adds
1{,}536 parameters to a 5.54\,M-parameter model, a relative increase of
$2.8 \times 10^{-4}$. Any measured difference between arms is therefore
attributable to the tokens' function, not to added capacity in the usual sense.

\subsection{Diagnostic measures}

\subsubsection{Layer-wise attention entropy}
For block $l$ with attention matrix
$A^{(l)} \in \mathbb{R}^{B \times H \times S \times S}$, whose rows are
probability distributions over keys, we compute the Shannon entropy of each row
and average over batch, heads and query positions:

\begin{equation}
\bar{H}^{(l)} = -\frac{1}{BHS}\sum_{b,h,i}\sum_{j=1}^{S} A^{(l)}_{b,h,i,j}\log_2\!\big(A^{(l)}_{b,h,i,j} + \epsilon\big),
\end{equation}

with $\epsilon = 10^{-12}$. Units are bits. A concentrated, artifact-like
routing pattern lowers $\bar{H}^{(l)}$; a diffuse one raises it. Because
recent \texttt{timm} builds execute fused scaled-dot-product attention, the
$[B,H,S,S]$ matrix is never materialised in the forward pass; our hook
reconstructs the softmax weights from the block's own $Q$ and $K$ projections,
reapplying the same scaling and any $q$/$k$ normalisation, so the recovered
matrix is the one the model used.

\subsubsection{Patch-norm outlier rate}
Following the artifact definition of \cite{darcet2024registers}, we take the
block-output residual stream $h^{(l)} \in \mathbb{R}^{B \times S \times d}$,
restrict it to the spatial positions $1+K \dots S$, and compute per-sample the
fraction of patch tokens whose $\ell_2$ norm exceeds three standard deviations
above the per-sample mean:

\begin{equation}
\rho^{(l)} = \frac{1}{BN}\sum_{b}\sum_{n} \mathbf{1}\!\left[\;\lVert h^{(l)}_{b,\,K+n}\rVert_2 > \mu^{(l)}_b + 3\sigma^{(l)}_b \;\right].
\end{equation}

The offset $1+K$ matters: slicing from index 1 would count register tokens as
image patches, and registers are expected to carry large norms by design, so
the error would manufacture exactly the artifact the metric is meant to detect.

\subsubsection{Generalization gap}
The regularization hypothesis is stated on losses rather than accuracies,
because loss is the quantity the optimiser acts on:

\begin{equation}
\Delta\mathcal{L}_s = \mathcal{L}^{\mathrm{val}}_{s}(e^{\star}_s) - \mathcal{L}^{\mathrm{train}}_{s}(E),
\end{equation}

where $e^{\star}_s$ is the epoch selected for seed $s$ and $E$ is the final
epoch. The gap is formed \emph{within} a seed before any averaging. Averaging
the two losses separately and subtracting would produce the same mean but
destroy the pairing, and with it any valid dispersion estimate for the gap;
since H1 is judged on that dispersion, the ordering of operations is part of
the method rather than an implementation detail.

\subsection{Statistical treatment}

With three seeds per arm, an unpaired comparison of arm means is dominated by
seed variance: the between-seed standard deviation of validation Top-1 in the
control arm ($2.17$ pp) is an order of magnitude larger than any treatment
effect we could plausibly detect. Every comparison in this paper is therefore
paired. For metric $m$, arm $K$ and the shared seed set
$\mathcal{S} = \{42, 1337, 3407\}$ we form

\begin{equation}
\delta_s = m_s(K) - m_s(0), \qquad s \in \mathcal{S},
\end{equation}

and report $\bar{\delta}$, its sample standard deviation $\hat{\sigma}_\delta$
(with $\mathrm{ddof} = 1$), the count of seeds moving in the hypothesised
direction, and the paired statistic
$t = \bar{\delta} / (\hat{\sigma}_\delta / \sqrt{3})$ against Student's $t$ with
$\mathrm{df} = 2$. Arm-level dispersion in Table~\ref{tab:main_results} uses the
population convention ($\mathrm{ddof} = 0$), which is recorded in the emitted
summary artifact so the estimator is never ambiguous.

Two cautions apply throughout and we state them once. First,
$\mathrm{df} = 2$ gives these tests very low power; a $p$-value here is
effect-size evidence, not a confirmatory test. Second, we report fifteen paired
comparisons without multiplicity correction, so an isolated small $p$ among
them is weak evidence. Accordingly we treat the direction count---how many of
three seeds move the same way---as the primary signal and the $p$-value as
secondary.
